% ---- score2 standard preamble (from: python3 score2_tools.py preamble) ----
\documentclass[10pt]{extarticle}
\usepackage[a4paper,margin=1.5cm]{geometry}
\usepackage[T1]{fontenc}
\usepackage{amsmath,amssymb}
\usepackage{graphicx}
\usepackage[export]{adjustbox}
\usepackage{multicol}
\usepackage{enumitem}
\usepackage{booktabs,array,tabularx}
\setlength{\columnsep}{0.7cm}
\setlength{\parindent}{0pt}
\setlength{\parskip}{3pt}
\pagestyle{empty}
\makeatletter
% figure registry (figdefs.tex is generated by the extract step)
\newcommand{\qfigdef}[4]{\@namedef{qf@file@#1}{#2}\@namedef{qf@w@#1}{#3}\@namedef{qf@wide@#1}{#4}}
\newcommand{\qf@check}[1]{\@ifundefined{qf@file@#1}{\PackageError{score2}{Unknown figure id #1}{}}{}}
\newcommand{\qf@img}[4][]{\includegraphics[width=#2\dimexpr\@nameuse{qf@w@#4}pt\relax,max width=#3,max height=0.45\textheight,#1]{\@nameuse{qf@file@#4}}}
% Full-width questions: finalize wraps every question that contains a wide figure in \qfullbegin...\qfullend,
% so the WHOLE question (stem, figures, options) is set across the page in reading order, then the two columns
% resume. The question's head - its number, stem and first figure (\qfullsplit marks the end of the head) -
% is kept on one page, so a number is never left at a page foot with its figure overleaf; the rest may break
% across pages normally, so no big gaps are left. Without a \qfullsplit, the whole question is kept together
% if it is under 45% of a page. Never write these yourself.
\newif\ifqf@full
\newif\ifqf@boxed
\newsavebox{\qf@box}
\newcommand{\qf@open}{\global\qf@fulltrue\par\global\qf@boxedtrue
  \setbox\qf@box\vbox\bgroup\hsize\textwidth\linewidth\textwidth}
\newcommand{\qf@emit}[1]{\par\egroup\global\qf@boxedfalse
  \ifdim\dimexpr\ht\qf@box+\dp\qf@box\relax<#1\textheight\noindent\usebox{\qf@box}\par\else\unvbox\qf@box\fi}
\newcommand{\qfullsplit}{\ifqf@boxed\qf@emit{0.85}\fi}
\newcommand{\qf@close}{\ifqf@boxed\qf@emit{0.45}\fi\global\qf@fullfalse}
\newcommand{\qfullbegin}{\end{multicols}\qf@open}
\newcommand{\qfullend}{\qf@close\begin{multicols}{2}}
\newcommand{\qfullstart}{\qf@open}
\newcommand{\qfulldone}{\qf@close}
% \qfig[scale]{id}: block figure, centred, directly where it is written (size set by the tool).
% Fallback if a wide figure is not inside a full-width question: break out of the columns in place - no floats.
\newcommand{\qfig}[2][1]{\qf@check{#2}%
  \ifqf@full\qf@blk{#1}{#2}\else
  \ifnum\@nameuse{qf@wide@#2}=1 \expandafter\@firstoftwo\else\expandafter\@secondoftwo\fi
  {\end{multicols}\begin{center}\qf@img{#1}{\textwidth}{#2}\end{center}\begin{multicols}{2}}%
  {\qf@blk{#1}{#2}}\fi}
\newcommand{\qf@blk}[2]{\par\smallskip{\centering\qf@img{#1}{\linewidth}{#2}\par}\smallskip}
% \qfigi[scale]{id}: inline figure (inside option lists / tables / minipages); never breaks columns.
\newcommand{\qfigi}[2][1]{\qf@check{#2}\qf@img[valign=c]{#1}{\linewidth}{#2}}
\makeatother
% structure macros - lint reads these to check completeness and figure placement
\newcommand{\papersec}[1]{\par\bigskip{\centering\Large\bfseries PAPER-#1\par}\medskip}
\newcommand{\partsec}[2]{\par\medskip\noindent{\large\bfseries PART-#1: #2}\par\smallskip}
\newcommand{\sect}[1]{\par\smallskip\noindent\textbf{#1}\par}
\newcommand{\qn}[1]{\par\smallskip\noindent\textbf{#1.}~}
\newcommand{\paraq}[2]{\par\smallskip\noindent\textbf{Paragraph for Questions #1 and #2}\par}
\newenvironment{qoptions}{\begin{itemize}[leftmargin=1.8em,itemsep=1pt,topsep=2pt,label={}]}{\end{itemize}}
\newcommand{\opt}[2]{\item[(#1)] #2}
% figure registry (generated by score2_tools.py finalize)
\qfigdef{p1_phy_s1_q05_1}{Score-I - Test 17 (17 Apr 2022)-img/p1_phy_s1_q05_1.png}{103.8}{0}
\qfigdef{p1_phy_s1_q06_1}{Score-I - Test 17 (17 Apr 2022)-img/p1_phy_s1_q06_1.png}{86.7}{0}
\qfigdef{p1_phy_s2_q10_1}{Score-I - Test 17 (17 Apr 2022)-img/p1_phy_s2_q10_1.png}{111.1}{0}
\qfigdef{p1_phy_s2_q11_1}{Score-I - Test 17 (17 Apr 2022)-img/p1_phy_s2_q11_1.png}{124.1}{0}
\qfigdef{p1_phy_s2_q12_1}{Score-I - Test 17 (17 Apr 2022)-img/p1_phy_s2_q12_1.png}{74.7}{0}
\qfigdef{p1_phy_s3_q05_1}{Score-I - Test 17 (17 Apr 2022)-img/p1_phy_s3_q05_1.png}{149.4}{0}
\qfigdef{p1_chem_s1_q02_1}{Score-I - Test 17 (17 Apr 2022)-img/p1_chem_s1_q02_1.png}{171.4}{0}
\qfigdef{p1_chem_s1_q05_B}{Score-I - Test 17 (17 Apr 2022)-img/p1_chem_s1_q05_B.png}{40.7}{0}
\qfigdef{p1_chem_s2_q09_1}{Score-I - Test 17 (17 Apr 2022)-img/p1_chem_s2_q09_1.png}{82.0}{0}
\qfigdef{p1_chem_s2_q11_A}{Score-I - Test 17 (17 Apr 2022)-img/p1_chem_s2_q11_A.png}{120.4}{0}
\qfigdef{p1_chem_s2_q11_B}{Score-I - Test 17 (17 Apr 2022)-img/p1_chem_s2_q11_B.png}{227.8}{0}
\qfigdef{p1_chem_s2_q11_C}{Score-I - Test 17 (17 Apr 2022)-img/p1_chem_s2_q11_C.png}{182.2}{0}
\qfigdef{p1_chem_s2_q11_D}{Score-I - Test 17 (17 Apr 2022)-img/p1_chem_s2_q11_D.png}{143.0}{0}
\qfigdef{p1_chem_s2_q12_B}{Score-I - Test 17 (17 Apr 2022)-img/p1_chem_s2_q12_B.png}{177.8}{0}
\qfigdef{p1_chem_s2_q12_C}{Score-I - Test 17 (17 Apr 2022)-img/p1_chem_s2_q12_C.png}{135.2}{0}
\qfigdef{p1_chem_s2_q12_D}{Score-I - Test 17 (17 Apr 2022)-img/p1_chem_s2_q12_D.png}{131.0}{0}
\qfigdef{p1_chem_s3_q05_1}{Score-I - Test 17 (17 Apr 2022)-img/p1_chem_s3_q05_1.png}{150.8}{0}
\qfigdef{p1_chem_s3_q06_1}{Score-I - Test 17 (17 Apr 2022)-img/p1_chem_s3_q06_1.png}{103.1}{0}
\qfigdef{p2_phy_s1_q01_1}{Score-I - Test 17 (17 Apr 2022)-img/p2_phy_s1_q01_1.png}{158.4}{0}
\qfigdef{p2_phy_s1_q03_1}{Score-I - Test 17 (17 Apr 2022)-img/p2_phy_s1_q03_1.png}{115.4}{0}
\qfigdef{p2_phy_s1_q04_1}{Score-I - Test 17 (17 Apr 2022)-img/p2_phy_s1_q04_1.png}{75.9}{0}
\qfigdef{p2_phy_s1_q06_1}{Score-I - Test 17 (17 Apr 2022)-img/p2_phy_s1_q06_1.png}{72.1}{0}
\qfigdef{p2_phy_s2_q01_1}{Score-I - Test 17 (17 Apr 2022)-img/p2_phy_s2_q01_1.png}{151.3}{0}
\qfigdef{p2_phy_s2_q02_1}{Score-I - Test 17 (17 Apr 2022)-img/p2_phy_s2_q02_1.png}{150.9}{0}
\qfigdef{p2_phy_s3_q03_1}{Score-I - Test 17 (17 Apr 2022)-img/p2_phy_s3_q03_1.png}{44.1}{0}
\qfigdef{p2_phy_s3_q04_1}{Score-I - Test 17 (17 Apr 2022)-img/p2_phy_s3_q04_1.png}{167.7}{0}
\qfigdef{p2_phy_s3_q05_1}{Score-I - Test 17 (17 Apr 2022)-img/p2_phy_s3_q05_1.png}{105.6}{0}
\qfigdef{p2_phy_s3_q06_1}{Score-I - Test 17 (17 Apr 2022)-img/p2_phy_s3_q06_1.png}{146.5}{0}
\qfigdef{p2_chem_s1_q05_1}{Score-I - Test 17 (17 Apr 2022)-img/p2_chem_s1_q05_1.png}{245.2}{0}
\qfigdef{p2_chem_s1_q05_3}{Score-I - Test 17 (17 Apr 2022)-img/p2_chem_s1_q05_3.png}{49.9}{0}
\qfigdef{p2_chem_s1_q05_2}{Score-I - Test 17 (17 Apr 2022)-img/p2_chem_s1_q05_2.png}{48.9}{0}
\qfigdef{p2_chem_s1_q05_4}{Score-I - Test 17 (17 Apr 2022)-img/p2_chem_s1_q05_4.png}{25.0}{0}
\qfigdef{p2_chem_s2_q01_1}{Score-I - Test 17 (17 Apr 2022)-img/p2_chem_s2_q01_1.png}{103.0}{0}
\qfigdef{p2_chem_s2_q05_1}{Score-I - Test 17 (17 Apr 2022)-img/p2_chem_s2_q05_1.png}{187.3}{0}
\qfigdef{p2_chem_s2_q06_1}{Score-I - Test 17 (17 Apr 2022)-img/p2_chem_s2_q06_1.png}{211.1}{0}
\qfigdef{p2_chem_s3_q05_1}{Score-I - Test 17 (17 Apr 2022)-img/p2_chem_s3_q05_1.png}{299.5}{1}
\qfigdef{p2_chem_s3_q06_1}{Score-I - Test 17 (17 Apr 2022)-img/p2_chem_s3_q06_1.png}{101.5}{0}
\qfigdef{p2_math_s2_q03_1}{Score-I - Test 17 (17 Apr 2022)-img/p2_math_s2_q03_1.png}{53.2}{0}
% ---- body skeleton ----
% \begin{document}
% \begin{center}{\Large\bfseries JEE (Main+Advanced) Score-II --- Test NN (DD Mon YYYY)}\end{center}
% \begin{multicols}{2}
%   \papersec{1}   (only for two-paper tests)
%   \partsec{1}{PHYSICS} \sect{Section-I (i): Single Correct} \qn{1} ... \qfig{p1_phy_q01_1} ...
% \end{multicols}
% \end{document}

\begin{document}
\begin{center}{\Large\bfseries JEE (Main+Advanced) Score-I --- Test 17 (17 Apr 2022)}\end{center}
\raggedcolumns
\setlength{\tabcolsep}{3pt}
\begin{multicols}{2}
\papersec{1}
\partsec{1}{PHYSICS}
\sect{SECTION-I(i)}
\qn{1} Atoms of a hydrogen like gas are in a particular excited energy level. When these atoms de-excite, they emit photons of different energies. Maximum and minimum energies of emitted photons are $E_{\max}=52.224$ eV and $E_{\min}=1.224$ eV respectively. Calculate principal quantum number of initially excited energy level. (Ionisation energy of hydrogen atom is 13.6 eV)
\begin{qoptions}
\opt{A}{6}
\opt{B}{4}
\opt{C}{5}
\opt{D}{3}
\end{qoptions}

\qn{2} An infinitely long metal cylinder of radius $R$ and surface charge density $\sigma$ is placed symmetrically with an imaginary surface of the shape of a prism. The length of prism is $L$ and its 3 sides are all equal to $3R$. The flux through the prism is :
\begin{qoptions}
\opt{A}{$\dfrac\pi3\dfrac{RL\sigma}{\varepsilon_0}$}
\opt{B}{$\dfrac\pi2\dfrac{RL\sigma}{\varepsilon_0}$}
\opt{C}{$\dfrac{\pi RL\sigma}{\varepsilon_0}$}
\opt{D}{$\dfrac{2\pi RL\sigma}{\varepsilon_0}$}
\end{qoptions}

\qn{3} A rough ball is thrown such that it rotates clockwise and moves right in viscous air. The force due to pressure difference of air acts :
\begin{qoptions}
\opt{A}{Upwards}
\opt{B}{Downwards}
\opt{C}{Nowhere}
\opt{D}{Can't say}
\end{qoptions}

\qn{4} Mercury is filled in a glass beaker which is heated by a flame from below. The level of mercury with time
\begin{qoptions}
\opt{A}{rises}
\opt{B}{falls}
\opt{C}{rises \& then falls}
\opt{D}{falls \& then rises}
\end{qoptions}

\qn{5} A bar magnet of mass $m$ is given an initial velocity $V_0$ towards the centre of a metallic coil of same mass initially at rest in gravity free space. The total loss of energy is :
\qfig{p1_phy_s1_q05_1}
\begin{qoptions}
\opt{A}{$\dfrac12mv_0^2$}
\opt{B}{$mv_0^2$}
\opt{C}{$\dfrac14mv_0^2$}
\opt{D}{cannot be determined}
\end{qoptions}

\qn{6} A string is wound on a spool which is kept on a rough horizontal floor having sufficient friction to prevent slipping. If string is pulled then
\qfig{p1_phy_s1_q06_1}
(1) String winds\\
(2) String unwinds\\
(3) Spool moves left\\
(4) Spool moves right\\
(5) Point $P$ moves left\\
(6) point $P$ moves right\\
Which of the above options are correct ?
\begin{qoptions}
\opt{A}{2, 3, 5}
\opt{B}{1, 4, 6}
\opt{C}{2, 4, 6}
\opt{D}{2, 3, 4}
\end{qoptions}

\sect{SECTION-I(ii)}
\qn{7} Two girls of equal mass $m$ jump off from the border line of a stationary carriage of mass $M$ with same horizontal velocity $u$ relative to the carriage in same direction. Neglect the effect of friction. Mark incorrect options.
\begin{qoptions}
\opt{A}{They will impart greater velocity to the carriage by jumping off simultaneously}
\opt{B}{They will impart greater velocity to the carriage by jumping one after the other}
\opt{C}{They will impart same velocity to the carriage in whatever manner they jump off}
\opt{D}{Insufficient data to reply}
\end{qoptions}

\qn{8} A body is suspended from a spring balance kept in a satellite. The reading of the balance is $W_1$ when the satellite goes in an orbit of radius $R$ and is $W_2$ when it goes in an orbit of radius $2R$. Mark incorrect options.
\begin{qoptions}
\opt{A}{$W_1=W_2$}
\opt{B}{$W_1<W_2$}
\opt{C}{$W_1>W_2$}
\opt{D}{$W_1\ne W_2$}
\end{qoptions}

\qn{9} Electric field intensity at a point $(x,y)$ is given by $\vec I=(12xy^3-4x)\,\hat i+18x^2y^2\,\hat j$ where all parameters are in S.I. units. Then choose the correct statement(s) out of the following :
\begin{qoptions}
\opt{A}{It represents a conservative electric field.}
\opt{B}{It represents a non-conservative electric field.}
\opt{C}{Electric potential at point $(1\,\text{m},1\,\text{m})$ is $-4$V if electric potential at origin is zero.}
\opt{D}{Electric potential at point $(1\,\text{m},1\,\text{m})$ cannot be defined.}
\end{qoptions}

\qn{10} For the circuit shown in figure, which of the following statements are correct ?
\qfig{p1_phy_s2_q10_1}
\begin{qoptions}
\opt{A}{Time constant of the circuit is 0.5 sec.}
\opt{B}{In steady state current through battery is 1.0A.}
\opt{C}{In steady state current through inductor is 1.0A.}
\opt{D}{In steady state potential difference across $2\,\Omega$ resistor is 4V.}
\end{qoptions}

\qn{11} A solid cylinder of radius $R$ is projected with velocity $v_0$ on a rough horizontal surface of coefficient of friction `$\mu$' at $t=0$ and at the same instant it has an angular velocity of $\dfrac{4v_0}R$ as shown in figure. Then
\qfig{p1_phy_s2_q11_1}
\begin{qoptions}
\opt{A}{The cylinder stops at $t=\dfrac{v_0}{\mu g}$ and then turns back.}
\opt{B}{At $t=\dfrac{5v_0}{3\mu g}$ the cylinder starts pure rolling.}
\opt{C}{At $t=\dfrac{v_0}{\mu g}$ angular velocity of the cylinder is $\dfrac{2v_0}R$.}
\opt{D}{At $t=\dfrac{3v_0}{\mu g}$ speed of the cylinder is $v_0$.}
\end{qoptions}

\qn{12} In a cylindrical region of radius $R$, there exists a time varying magnetic field $B$ such that $\dfrac{dB}{dt}=k\,(>0)$. A charged particle having charge $q$ is placed at the point $P$ at a distance $d\,(>R)$ from its centre $O$. Now, the particle is moved in the direction perpendicular to $OP$ (see figure) by an external agent upto infinity so that there is no gain in kinetic energy of the charged particle. Choose the correct statement/s.
\qfig{p1_phy_s2_q12_1}
\begin{qoptions}
\opt{A}{Work done by external agent is $\dfrac{q\pi R^2}4k$ if $d=2R$}
\opt{B}{Work done by external agent is $\dfrac{q\pi R^2}8k$ if $d=4R$}
\opt{C}{Work done by external agent is $\dfrac{q\pi R^2}4k$ if $d=4R$}
\opt{D}{Work done by external agent is $\dfrac{q\pi R^2}4k$ if $d=6R$}
\end{qoptions}

\sect{SECTION-II}
\qn{1} A force $\vec F=-K\left(y\,\hat i+x\,\hat j\right)$, (where $K$ is a $+$ve constant) acts on a particle moving in $xy$ plane starting from origin, the particle is taken along the positive $x$-axis to the point $(a,0)$ and then parallel to $y$ axis to the point $(a,a)$. If the total work done by force $\vec F$ on the particle is $nka^2$, find $n+3$.

\qn{2} A mercury thermometer reads 80\textdegree C when the mercury is at 5.2 cm mark and 60\textdegree C when the mercury is at 3.9 cm mark. If the temperature when the mercury level is at 2.6 cm mark is $T$ (in degree Celsius), find $T-35$

\qn{3} A standing wave with a frequency of 1000 hz in a column of mixture of unknown gases at 27\textdegree C produces node that are 0.1 m apart. Find the value of $r=\dfrac{C_P}{C_V}$ for the mixture. Use $R=8.3$ J/mol--K and molar mass of mixture 124.5 gm/mol. $C_P$ \& $C_V$ are the molar heat capacities at constant pressure and constant volume respectively.

\qn{4} When a source of sound of frequency 600 Hz moves with velocity 30 m/s towards a stationary observer, the apparent wavelength is $\lambda$. If the same source is kept stationary and the observer is moved with same velocity 30 m/s towards the source, then apparent wavelength is $\lambda'$, then find $\lambda'-\lambda$ (in cm) for speed of sound in air to be 340 m/s.

\qn{5} Two rigid rods of same material and of length $L$ and $\dfrac L2$ are placed in equilibrium on a smooth horizontal plane at temperature $T$ as shown in figure. Springs are in their natural length. If the temperature is increased by $\Delta T$ amount energy stored in springs is $\dfrac9{5\beta}k\alpha^2L^2\Delta T^2$, where $k$ is the spring constant of spring ($A$) and $\alpha$ is the coefficient of linear expansion of material. (Note: Neglect the change in spring constant due to increase in temperature).
\qfig{p1_phy_s3_q05_1}

\qn{6} A particle is moving in a plane according to the law $x=at^2$, $y=bt^3$ where $a=2$ m/s$^2$, $b=1$ m/s$^3$. Radius of curvature of particle at $t=1$ s is found to be $\dfrac{\eta^3}{12}$ m. Find $\eta$.

\partsec{2}{CHEMISTRY}
\sect{SECTION-I(i)}
\qn{1} Which of the statement is incorrect?
\begin{qoptions}
\opt{A}{Hyponitrous acid is a weak acid and it is isomeric with nitramide.}
\opt{B}{Sodium dihydrogen pyrophosphate does not exist.}
\opt{C}{[Co(en)$_2$SO$_4$]$^{\oplus}$ is a low spin optically active compound.}
\opt{D}{Acidic character (w.r.t. oxide acceptor) : XeF$_6$ $>$ XeO$_2$F$_4$ $>$ XeO$_3$}
\end{qoptions}

\qn{2} Born Haber cycle helps to determine lattice enthalpy of ionic compound as well as electron gain enthalpy of different elements. Based on the diagram drawn, answer the questions given below.
\qfig{p1_chem_s1_q02_1}
What is the magnitude of lattice enthalpy of MgCl$_2$ ?
\begin{qoptions}
\opt{A}{20 kJ}
\opt{B}{40 kJ}
\opt{C}{15 kJ}
\opt{D}{40 kJ}
\end{qoptions}

\qn{3} Which of the following statement(s) is(are) correct?

(1) When a chalk stick is dipped in ink, both absorption and adsorption takes place and surface of chalk retains the color of ink while the inside in white.

(2) Water vapours are absorbed by silica gel

(3) Water vapours are adsorbed by anhydrous CaCl$_2$

(4) The precipitate of Mg(OH)$_2$ attains pink colour when precipitated in presence of magneson reagent
\begin{qoptions}
\opt{A}{4 only}
\opt{B}{1 Only}
\opt{C}{2, 3 only}
\opt{D}{1, 2, 3 only}
\end{qoptions}

\qn{4} Density of a Vander Waal's gas at 500 K and 1 atm is found to be 0.8 kg/m$^3$. It is also found that the gas diffuses $\sqrt{\dfrac32}$ times slower than pure O$_2$ gas under identical conditions. The value of `z' for the real gas is (R = 0.08 atm-$\ell$/mol-K) :
\begin{qoptions}
\opt{A}{2/3}
\opt{B}{3/2}
\opt{C}{1.8}
\opt{D}{0.8}
\end{qoptions}

\qn{5} Which of the following oxyanion has 3 identical resonating structures?
\begin{qoptions}
\opt{A}{CO$_3^{2-}$}
\opt{B}{\qfigi{p1_chem_s1_q05_B}}
\opt{C}{ClO$_4^{-}$}
\opt{D}{PO$_4^{3-}$}
\end{qoptions}

\qn{6} (+)2--butanol has $[\theta]_\lambda^{25}=13.9^\circ$. A sample of 2-butanol containing both the enantiomers was found to have specific rotation value of $-3.5^\circ$ under similar condition. The percentage of the (+) and ($-$) enantiomer present in the sample are, respectively :
\begin{qoptions}
\opt{A}{37.4\% and 62.6\%}
\opt{B}{62.6\% and 37.4\%}
\opt{C}{42.2\% and 57.8\%}
\opt{D}{35.5\% and 64.5\%}
\end{qoptions}

\sect{SECTION-I(ii)}
\qn{7} Which of the following do \textbf{NOT} represent the correct sequence of indicated properties ?
\begin{qoptions}
\opt{A}{[FeF$_6$]$^{3-}$ $>$ [NiF$_6$]$^{4-}$ : Magnetic moment}
\opt{B}{[Fe(CN)$_6$]$^{3-}$ $<$ [Fe(NH$_3$)$_6$]$^{3+}$ : stability constant}
\opt{C}{NiO $>$ CaO $>$ BeO : Basic character}
\opt{D}{F$_2$ $<$ Cl$_2$ $<$ Br$_2$ $<$ I$_2$ : Colour intensity}
\end{qoptions}

\qn{8} Which of the following is/are true statements?
\begin{qoptions}
\opt{A}{In Puddling process wrought iron is produced.}
\opt{B}{In blast furnace phosphate and silicate are reduced into P and Si which pass into molten iron.}
\opt{C}{B, Ga, In, Ge and Si are refined by zone refining.}
\opt{D}{In Mond's process high spin complex is formed.}
\end{qoptions}

\qn{9} If an optically active compound `A' decompose through given parallel 1$^{\text{st}}$ order kinetics.
\qfig{p1_chem_s2_q09_1}
Initial mole of A is 2 ( $k_1=0.0693\ \text{sec}^{-1}$, $k_2=0.1386\ \text{sec}^{-1}$). If only A, B \& D are optically active compounds and their angle of rotation per mole are $60^\circ$, $30^\circ$, $-90^\circ$ respectively, then which of the following is/are correct -
\begin{qoptions}
\opt{A}{Angle of rotation after time 3.33 sec is $20^\circ$}
\opt{B}{Mixture is levorotatary after 75\% decomposition of A}
\opt{C}{Overall $t_{1/2}$ of A is 10/3 sec.}
\opt{D}{Percentage of C formed in the product is 11.11\%}
\end{qoptions}

\qn{10} 100 ml, 0.3 M solution of HCl is mixed with 50 of water to form solution (A). 50 ml of solution (A) is mixed with 50 ml, 0.05M NaOH to form solution (B). Another 50 ml of solution (A) is mixed with 100 ml, 0.05M NH$_4$OH to form solution (C). Last 50 ml of solution (A) is mixed with 100 ml, 0.1 M of AgNO$_3$ to form solution (D),

[Given K$_{\text{sp}}$(AgCl) = $10^{-10}$, K$_{\text{b}}$ (NH$_4$OH) = $2\times10^{-5}$, log 2 = 0.3, log 3 = 0.48)]
\begin{qoptions}
\opt{A}{pH of solution (A) is 0.7}
\opt{B}{pH of solution (B) is 1.12}
\opt{C}{pH of solution (C) is 4.7}
\opt{D}{mole of Ag$^+$ ion in solution D is $1.5\times10^{-6}$ mole}
\end{qoptions}

\qn{11} Which of the following reactions give para bromophenol as a major product?
\begin{qoptions}
\opt{A}{\qfigi{p1_chem_s2_q11_A}}
\opt{B}{\qfigi{p1_chem_s2_q11_B}}
\opt{C}{\qfigi{p1_chem_s2_q11_C}}
\opt{D}{\qfigi{p1_chem_s2_q11_D}}
\end{qoptions}

\qn{12} In which of the following reaction(s) the product is/are correctly matched?
\begin{qoptions}
\opt{A}{CH$_3$--CH=O $\xrightarrow[\text{(ii) H}^+]{\text{(i) [Ag(NH}_3)_2]^+}$ CH$_3$--COOH}
\opt{B}{\qfigi{p1_chem_s2_q12_B}}
\opt{C}{\qfigi{p1_chem_s2_q12_C}}
\opt{D}{\qfigi{p1_chem_s2_q12_D}}
\end{qoptions}

\sect{SECTION-II}
\qn{1} Among the following cations(s) :

Hg$^{2+}$, Mg$^{2+}$, Cu$^{2+}$, Ag$^{+}$, Pb$^{2+}$

(i) Number of cation(s) gives(s) ppt with H$_2$S / H$^+$ is = (X)

(ii) Number of cation(s) gives(s) ppt with KI is = (Y)

(iii) Number of cation(s) gives(s) ppt with dil. HCl is = (Z)

Find the value of (X + Y + Z)

\qn{2} Consider the following reactions\\
(i) PbS + H$_2$O$_2$ $\to$\\
(ii) Ag$_2$O + H$_2$O$_2$ $\to$\\
(iii) K$_4$[Fe(CN)$_6$] + H$_2$O$_2$ + H$_2$SO$_4$ $\to$\\
(iv) K$_3$[Fe(CN)$_6$] + H$_2$O$_2$ + KOH $\to$\\
(v) K$_2$Cr$_2$O$_7$ + H$_2$O$_2$ + H$_2$SO$_4$ $\xrightarrow[\text{heating}]{\Delta}$\\
(vi) CrCl$_3$ + H$_2$O$_2$ + KOH $\to$\\
(vii) Cl$_2$ + H$_2$O$_2$ $\xrightarrow{\text{H}^+}$\\
(viii) KIO$_4$ + H$_2$O$_2$ $\xrightarrow{\text{H}^+}$\\
(ix) NaMnO$_4$ + H$_2$O$_2$ + H$_2$SO$_4$ $\to$

In how many reactions oxidising nature of H$_2$O$_2$ is observed.

\qn{3} A definite amount of an ideal gas is taken from state-I to state-II isothermally at 300K. If the process is carried out reversibly then work done is $-$ 1750 kJ and if process is carried out irreversibly then work done is $-$1000 kJ. If $\Delta S_{\text{total}}$ is x kJ/K for irreversible isothermal process then give `x' .

\qn{4} Liquids A and B form ideal solution over the entire range of composition. At temperature T, equimolar binary solution of liquids A and B has vapour pressure 45 Torr. At the same temperature, a new solution of A and B having mole fractions $x_A$ and $x_B$ respectively, has vapour pressure of 22.5 Torr. If the value of $\dfrac{x_A}{x_B}$ in the new solution is Y. Then the value `Y' is :

(Given that the vapour pressure of pure liquid A is 20 Torr at temperature T)

\qn{5} Consider the figure
\qfig{p1_chem_s3_q05_1}
\begin{center}\footnotesize
\begin{tabular}{|c|l|c|c|c|}
\hline
S.no. & Amino acids & \multicolumn{3}{c|}{pKa values}\\
\cline{3-5}
 & & $\alpha$-COOH & $\alpha$-N$^+$H$_3$ & Side chain\\
\hline
1 & Arginine & 2.17 & 9.04 & 12.48\\
\hline
2 & Alanine & 2.34 & 9.69 & ........\\
\hline
3 & Aspartic acid & 2.09 & 9.82 & 3.86\\
\hline
4 & Lysine & 2.18 & 8.95 & 10.79\\
\hline
5 & Glutamic acid & 2.19 & 9.67 & 4.25\\
\hline
6 & Valine & 2.32 & 9.62 & ........\\
\hline
\end{tabular}
\end{center}
Calculate :\\
Number of amino acids move towards cathode = X\\
Number of non essential amino acids move towards anode = Y\\
Find $\dfrac{X}{Y}$ = ?

\qn{6} How many reagents can be used to distinguish these compounds :
\qfig{p1_chem_s3_q06_1}
(i) NaHCO$_3$ \quad (ii) Na metal \quad (iii) Neutral FeCl$_3$\\
(iv) AgNO$_3$ + NH$_4$OH \quad (v) CS$_2$, HgCl$_2$ \quad (vi) NaOI\\
(vii) 2,4--DNP

\partsec{3}{MATHEMATICS}
\sect{SECTION-I(i)}
\qn{1} If $z\displaystyle\int_0^z \sin(f(x))\,dx=(z+2)\int_0^z x\sin(f(x))\,dx$ where $z>0$, then $\left|(1+z)f'(z)\cot(f(z))\right|$ is equal to
\begin{qoptions}
\opt{A}{0}
\opt{B}{1}
\opt{C}{2}
\opt{D}{3}
\end{qoptions}

\qn{2} Given vectors $\vec a=\sin\dfrac\pi6\,\hat i+\sin\dfrac\pi3\,\hat j$; $\vec b=\sin\dfrac\pi3\,\hat j+\sin\dfrac\pi6\,\hat k$ and $\vec c=\hat i+2\hat j+3\hat k$. The volume of parallelopiped with three conterminous edges as

$\vec u=(\vec a\cdot\vec a)\,\vec a+(\vec a\cdot\vec b)\,\vec b+(\vec a\cdot\vec c)\,\vec c$

$\vec v=(\vec a\cdot\vec b)\,\vec a+(\vec b\cdot\vec b)\,\vec b+(\vec b\cdot\vec c)\,\vec c$

$\vec w=(\vec a\cdot\vec c)\,\vec a+(\vec b\cdot\vec c)\,\vec b+(\vec c\cdot\vec c)\,\vec c$ equals
\begin{qoptions}
\opt{A}{$\left(2\cos\dfrac\pi4-\sin\dfrac\pi3\right)^3$}
\opt{B}{$\left(2\cos\dfrac\pi3-\sin\dfrac\pi3\right)^3$}
\opt{C}{$\left(2\cos\dfrac\pi6-\sin\dfrac\pi6\right)^3$}
\opt{D}{$\left(2\sin\dfrac\pi4-\cos\dfrac\pi3\right)^3$}
\end{qoptions}

\qn{3} A parabola whose focus S(3, 4) is touching the co-ordinate axes. The equation of circle whose diameter is the portion of tangent at vertex of the parabola between the co-ordinate axes is
\begin{qoptions}
\opt{A}{$x^2+y^2-6x-8y=0$}
\opt{B}{$x^2+y^2+3x+4y=0$}
\opt{C}{$x^2+y^2+6x+8y=0$}
\opt{D}{$x^2+y^2-3x-4y=0$}
\end{qoptions}

\qn{4} Number of complex number(s) `z' satisfying $|z|^n=|z|^{n-2}z^2+|z|^{n-2}z+1$ ($n\in\mathbb Q$) and $\text{Re}(z)\neq\dfrac{-1}{2}$ is(are)
\begin{qoptions}
\opt{A}{0}
\opt{B}{1}
\opt{C}{2}
\opt{D}{4}
\end{qoptions}

\qn{5} A polynomial function P(x) of degree 5 with leading coefficient one, increases in interval $(-\infty,1)$ and $(3,\infty)$ and decreases in interval $(1,3)$. Given that $P'(2)=0$ then value of $P'(4)$ is
\begin{qoptions}
\opt{A}{240}
\opt{B}{180}
\opt{C}{120}
\opt{D}{60}
\end{qoptions}

\qn{6} Let $f(x)=\begin{vmatrix}\sec x&\cos x&\sec^2x+\cot x\csc x\\ \cos^2x&\cos^2x&\csc^2x\\ 1&\cos^2x&\cos^2x\end{vmatrix}$, then $\displaystyle\int_{-\pi}^{\pi}f(x)\,dx$ has the value equal to
\begin{qoptions}
\opt{A}{$\dfrac{-\pi}{4}$}
\opt{B}{$\dfrac{-\pi}{2}$}
\opt{C}{$\pi$}
\opt{D}{$-\pi$}
\end{qoptions}

\sect{SECTION-I(ii)}
\qn{7} If $z$ is non-real complex number satisfying $|z|=2$, then which of the following is(are) correct ?
\begin{qoptions}
\opt{A}{$\arg\!\left(\dfrac{z-2}{z+2}\right)=\pm\dfrac\pi2$ (where $z\neq\pm2$)}
\opt{B}{$\arg\!\left(\dfrac{z+1+i\sqrt3}{z-1+i\sqrt3}\right)=\dfrac\pi6$}
\opt{C}{$|z^2-1|\ge3$}
\opt{D}{$|z^2-1|\le5$}
\end{qoptions}

\qn{8} Let $\displaystyle\sum_{i=0}^{100}\sum_{j=0}^{100}\left(C_i^{\,2}+C_j^{\,2}+C_iC_j\right)=\text{a}\cdot{}^{2n}C_n+2^b$, where a, b, n are even natural numbers and $C_r$ represent the coefficient of $x^r$ in the expansion of $(1+x)^{100}$ then
\begin{qoptions}
\opt{A}{a + b + n = 504}
\opt{B}{a + b = 402}
\opt{C}{a $-$ n = 102}
\opt{D}{a + 2b + 3n = 900}
\end{qoptions}

\qn{9} Let $f(x)=2\tan^{-1}x$ and $g(x)$ be a differentiable function satisfying $g\!\left(\dfrac{x+2y}{3}\right)=\dfrac{g(x)+2g(y)}{3}$ $\forall\ x,y\in R$ and $g'(0)=1$ and $g(0)=2$. The number of integers `x' satisfying $f^2(g(x))-5f(g(x))+4>0$ where $x\in(-10,20)$ is greater than or equal to
\begin{qoptions}
\opt{A}{5}
\opt{B}{6}
\opt{C}{7}
\opt{D}{8}
\end{qoptions}

\qn{10} Let six children $B_1,B_2,B_3,B_4,B_5,B_6$ each have a fair 6 sided die which they roll simultaneously. Let E denotes the event that atleast three of them roll a prime number and F denotes the event that $B_1,B_2,B_3$ throw a prime number. If the conditional probability of F given E is equal to $\dfrac ab$ where a and b are coprime then
\begin{qoptions}
\opt{A}{$2a+b=29$}
\opt{B}{$|a-b|=10$}
\opt{C}{$ab\le84$}
\opt{D}{$b-4a=0$}
\end{qoptions}

\qn{11} If pair of variable straight lines $x^2+4y^2+\alpha xy=0$ ($\alpha$ is real parameter) cuts the ellipse $x^2+4y^2=4$ at two points A and B, then which of the following cannot be the locus of point of intersection of tangents at A and B of ellipse
\begin{qoptions}
\opt{A}{$x^2+4y^2=0$}
\opt{B}{$4x^2-y^2=0$}
\opt{C}{$x^2-xy+4y^2$}
\opt{D}{$x^2-xy-4y^2=0$}
\end{qoptions}

\qn{12} A non-zero vector $\vec p$ is equally inclined to three vectors $\vec a=\hat i-\hat j+\hat k$, $\vec b=2\hat i+\hat j$, $\vec c=3\hat j-2\hat k$. Let $\vec x$, $\vec y$, $\vec z$ be three vectors in plane of $\vec a$; $\vec b$; $\vec b$; $\vec c$ and $\vec c$; $\vec a$ respectively then
\begin{qoptions}
\opt{A}{$\vec x\cdot\vec p=1$}
\opt{B}{$\vec y\cdot\vec p=0$}
\opt{C}{$\vec z\cdot\vec p=0$}
\opt{D}{$\left(2\hat i-3\hat j+4\hat k\right)\cdot\vec p=0$}
\end{qoptions}

\sect{SECTION-II}
\qn{1} Let $f(x)=x^3-ax^2+bx+6$ where $a,b\in\{1,2,3,4,5,6\}$. The probability that $f(x)$ is strictly increasing is

\qn{2} Let $a_1,a_2,a_3,\ldots,a_n$ be an increasing arithmetic progression of positive integers. If $a_3=13$, then the maximum value of $\displaystyle\sum_{n=1}^{5}a_{a_n}$ is

\qn{3} Let `C' be the curve passing through the point (1, 1) has the property that the perpendicular distance of the normal from origin at any point P of the curve is equal to distance of P from x-axis. If area bounded by curve `C' and x-axis in first quadrant is $k\pi$ square units then k is equal to

\qn{4} A circle is inscribed in a triangle with side lengths 3, 4, 5. A second circle interior to the triangle, is tangent to the first circle and to both sides of the larger acute angle of triangle. The radius of the second circle is

\qn{5} Let f be a function defined by $f(x)=\begin{cases}(x-r)^2 & ;\ r-1\le x\le r+1\\ 1 & ;\ r+1\le x\le r+2\end{cases}$, where $r=3k$, $k\in I$ then value of $\displaystyle\int_0^{30}f(x)\,dx$ is equal to

\qn{6} Let $A=[a_{ij}]$ be a $3\times3$ matrix such that


$a_{ij}=\displaystyle\lim_{x\to0}\left((\sin ix)^{j/x}+\dfrac{1-\cos jx}{ix}\right)$ ; $i=j$

$a_{ij}=\sin^{2021}\!\left(i^{2020}-j^{2020}\right)x+2\tan^{-1}\!\left(e^{(i-j)x}\right)-\dfrac\pi2$ ; $i\neq j$

then the value of $\det(\operatorname{adj}A)$ is

\papersec{2}
\partsec{1}{PHYSICS}
\sect{SECTION-I}
\qn{1} A vertical frictionless semi-circular track of radius 1 m is fixed on the edge of a movable trolley as shown in figure. Initially the system is rest and a mass $m$ is kept at the top of the track. The trolley starts moving to the right with a uniform horizontal acceleration $a=2g/9$. The mass slides down the track, eventually losing contact with it and dropping to the floor.
\qfig{p2_phy_s1_q01_1}
\begin{qoptions}
\opt{A}{The angle $\theta$ with vertical at which it loses contact with the trolley is $37^\circ$.}
\opt{B}{The angle $\theta$ with vertical at which it loses contact with the trolley is $53^\circ$.}
\opt{C}{The time taken by the mass to drop on the floor, after losing contact is 0.43 s.}
\opt{D}{The time taken by the mass to drop on the floor, after losing contact is 0.51 s.}
\end{qoptions}

\qn{2} Each egg of Caraphractus cinctus, a parasitic wasp, has a mass of about $2.0\times10^{-10}$ kg. Consider the formation of $^{16}_{8}$O from H atoms and neutrons. Let `n' be the number of nuclei of $^{16}_{8}$O that must be formed to produce a mass defect equal to the mass of one Caraphractus cinctus egg. The given atomic masses are: $M(^{16}_{8}\text{O})=15.994\,915$ u, $M(^{1}_{1}\text{H})=1.007\,825$ u and $m_n=1.008\,665$ u.
\begin{qoptions}
\opt{A}{$n=8.8\times10^{17}$}
\opt{B}{$n=8.8\times10^{16}$}
\opt{C}{The total binding energy of these $^{16}_{8}$O nuclei is $1.1\times10^{20}$ MeV}
\opt{D}{The total binding energy of these $^{16}_{8}$O nuclei is 125 MeV}
\end{qoptions}

\qn{3} A parallel-plate capacitor is charged as shown (Q is given). A metal slab with the total charge $+Q$ is placed inside the capacitor as shown. The thickness of the slab is $d$. The distance between the top plate and the top of the slab is $2d$, and the distance between the bottom plate and the bottom of the slab is $d$. Each plate is grounded through a galvanometer as shown. Find the charge that passes through each galvanometer after both switches are closed simultaneously.
\qfig{p2_phy_s1_q03_1}
\begin{qoptions}
\opt{A}{A total charge of $-Q/3$ from through G1 from the bottom plate to the ground.}
\opt{B}{A total charge of $Q/3$ from through G1 from the bottom plate to the ground.}
\opt{C}{A total charge of $4Q/3$ from through G2 from the top plate to the ground.}
\opt{D}{A total charge of $-4Q/3$ from through G2 from the top plate to the ground.}
\end{qoptions}

\qn{4} Two capacitors are connected by the ideal wires forming a square with side $d$. The central wire connects the midpoints of the opposite sides of the ``square'' as shown. The circuit is placed in the magnetic field, directed into the page, that changes with time as $B(t)=kt$, where $k$ is a constant coefficient. After some time, the central wire is cut, and the magnetic field is then kept constant. Find the charges on the capacitors after the equilibrium has been reached. The capacitances $C_1$ and $C_2$ are given.
\qfig{p2_phy_s1_q04_1}
\begin{qoptions}
\opt{A}{$Q_1=\left(\dfrac{C_2-C_1}{C_2+C_1}\right)\dfrac{C_2kd^2}{2}$}
\opt{B}{$Q_1=\left(\dfrac{C_2-C_1}{C_2+C_1}\right)\dfrac{C_1kd^2}{2}$}
\opt{C}{$Q_2=\left(\dfrac{C_2-C_1}{C_2+C_1}\right)\dfrac{C_1kd^2}{2}$}
\opt{D}{$Q_2=\left(\dfrac{C_2-C_1}{C_2+C_1}\right)\dfrac{C_2kd^2}{2}$}
\end{qoptions}

\qn{5} A thin film of magnesium fluoride ($\mu=1.38$) is used to coat a camera lens ($\mu=1.56$). Which of the following thickness of coating will not allow any strong reflection for the wavelength 552 nm for near normal incidence ?
\begin{qoptions}
\opt{A}{1000 nm}
\opt{B}{500 nm}
\opt{C}{300 nm}
\opt{D}{100 nm}
\end{qoptions}

\qn{6} Two small spheres of mass $m$ each are attached to the ends of a light rigid rod bent at a right angle and pivoted in the vertical plane at the vertex of the right angle. The vertical part of the rod is three times as long as the horizontal part. Initially, the rod is held at rest in the position shown in the diagram. The pivot is frictionless. Force exerted by the rod on the pivot immediately after the rod is released:
\qfig{p2_phy_s1_q06_1}
\begin{qoptions}
\opt{A}{in horizontal direction is released is $3mg/10$}
\opt{B}{in horizontal direction is released is $19mg/10$}
\opt{C}{in vertical direction is released is $3mg/10$}
\opt{D}{in vertical direction is released is $19mg/10$}
\end{qoptions}

\sect{SECTION-II}
\qn{1} A wire has been put in the shape of the curve $r=\dfrac{1}{\cos\theta\,(1+\sin^2\theta)}$ as shown in figure. The wire is infinite in length and carries a current $I=10$ A. Find the magnetic field (in $\mu$T) due to this current carrying wire at the origin. $r$ and $\theta$ are two dimensional polar co-ordinates in the plane of the paper.
\qfig{p2_phy_s2_q01_1}

\qn{2} A research vessel has a round glass window in the bottom for observing the seabed. The diameter of the window is 60 cm, the thickness of the glass is 20 mm, and the index of refraction of water is $\dfrac43$; that of the glass is 1.55. The seabed is 6.0 m beneath the window. Find the area (in m$^2$) of the seabed that can be seen through the window.
\qfig{p2_phy_s2_q02_1}

\qn{3} In 1994, a unicycle with a wheel diameter of 2.5 cm was ridden 3.6 m in Las Vegas, Nevada. Suppose the wheel has 200 turns of thin wire wrapped around its rim, creating loops with the same diameter as the wheel. An emf of 9.6 mV is induced when the wheel is perpendicular to a magnetic field that steadily decreases from 0.68 T to 0.24 T. For how long (in seconds) is the emf induced ?

\qn{4} A projectile launched from the ground explodes into three fragments of equal mass at the top point of the trajectory. One of the fragments lands $t=3$ seconds after the explosion; two other fragments land simultaneously $2t$ seconds after the explosion. How high (in meters) above the ground does the projectile explode?

\qn{5} An alien spaceship is revolving around a spherical asteroid in a very low circular orbit. The period of revolutions of the satellite is 15.0 minutes. A research party of the alien students lands on the asteroid and drills a tunnel through it along the diameter. They then launch a rock into the tunnel with the initial speed equal to that of the orbiting spaceship. How long (in minutes) would it be before the rock comes back to the launch point?

\qn{6} An open glass tube has a length of 0.500 m. The tube is held vertically above a water-filled container so that the bottom end of the tube slightly touches the surface of the water. Under these conditions, the height of the water column in the tube is 20.0 mm. What would be the height (in mm) of the water column in the tube if the top end of the tube is sealed before touching the water? The atmospheric pressure is $1.01\times10^5$ Pa.

\sect{SECTION-III}
\qn{1} A small, charged ball is hovering in the state of equilibrium at a height $h$ over a large horizontal uniformly charged dielectric plate. What would be the approximate acceleration of the ball (in $\mu$m/s$^2$) if a disc of radius $r=0.001h$ is removed from the plate directly underneath the ball?

\qn{2} One mole of a monatomic gas is heated in such a way that its molar specific heat is $2R$. During the heating, the volume of the gas is doubled. By what factor does the temperature of the gas change?

\qn{3} A container with helium is sealed by a movable heavy piston as shown. Initially, the piston is in equilibrium. The piston is then slowly lifted a distance $L=5$ cm and held until the thermal equilibrium with the surroundings is re-established. After that, the container is thermally insulated from the surroundings and piston is released. What approximately is the final equilibrium position of the piston with respect to its initial position (in cm)? Neglect the outside air pressure on the piston and the heat capacities of the piston and the container. Assume $V\gg AL$
\qfig{p2_phy_s3_q03_1}

\qn{4} Determine potential difference across capacitor at $t=\dfrac{\pi}{16}$ sec. [Emf of AC source is in volts and `t' is in seconds]
\qfig{p2_phy_s3_q04_1}

\qn{5} For the vernier calliper shown in the figure 1 msd = 1 mm and 7 msd = 8 vsd. The instrument has no zero error. The length of the rod is $4\left(\dfrac{x}{8}\right)$ mm. Find `x'.
\qfig{p2_phy_s3_q05_1}

\qn{6} In the diagram, a platform of mass $M$ is placed on a horizontal frictionless surface. Two blocks of masses $4M$ and $M$ are placed on the platform. For both blocks, the coefficients of static friction with the platform equal 0.16 and the coefficients of kinetic friction equal 0.10. The blocks are connected by a light string through a light pulley, which is acted upon by an unknown horizontal force $F$. If the acceleration of the platform is 2 m/s$^2$, find the acceleration (in m/s$^2$) of block of mass `M'.
\qfig{p2_phy_s3_q06_1}

\partsec{2}{CHEMISTRY}
\sect{SECTION-I}
\qn{1} One of the major use of H$_2$O$_2$ is for artificial respiration. It is stored in cylinder of 3 litre capacity completely filled with H$_2$O$_2$ solution of `45.4V'. In a hospital of kota, patient under artificial respiration takes 200 ml O$_2$ per min at 1 atm at 273 K for 2.8 hours from cylinder. [Assume volume of solution and rate of decomposition remain constant]

Then which of the following statement(s) is/are correct.
\begin{qoptions}
\opt{A}{Initial moles of H$_2$O$_2$ in a cylinder are 12}
\opt{B}{Volume of oxygen used for respiration is 33.6 L}
\opt{C}{Volume strength of H$_2$O$_2$ left in solution is 11.35 V}
\opt{D}{Volume strength of H$_2$O$_2$ left in solution is 34.05 V}
\end{qoptions}

\qn{2} Which of the following statements are correct -
\begin{qoptions}
\opt{A}{Amorphous solids are also called pseudo solids or supercooled liquids.}
\opt{B}{If a solid is very hard, and it is electrical insulator in solid as well as in molten state and melts at extremely high temperature, then it can be covalent or network solid}
\opt{C}{In Rhombohedral crystal system, only primitive unit cell is possible}
\opt{D}{All atoms of `B' layer in ABAB type of close packing in 3D is completely inside unit cell}
\end{qoptions}

\qn{3} CORRECT order/s of ionisation energy is/are :
\begin{qoptions}
\opt{A}{Pb$^+$ $>$ Sn$^+$}
\opt{B}{Na$^+$ $>$ Mg$^+$}
\opt{C}{Li$^+$ $<$ O$^+$}
\opt{D}{Be$^+$ $<$ C$^+$}
\end{qoptions}

\qn{4} Select the correct statement(s) :-
\begin{qoptions}
\opt{A}{Compounds [Cr(H$_2$O)$_6$]Cl$_3$ and [Cr(H$_2$O)$_4$Cl$_2$]\allowbreak Cl.2H$_2$O are hydrate isomers.}
\opt{B}{Greater the oxidation state of central metal cation, higher is the value of $\Delta$(CFSE).}
\opt{C}{The stability constant of [Cu(NH$_3$)$_4$]$^{2+}$ is more than that of [CuCl$_4$]$^{2-}$}
\opt{D}{In the complex ion [CoF$_6$]$^{3-}$, F$^-$ is a weak field ligand, so that $\Delta_o<$ P and it is a low spin complex.}
\end{qoptions}

\qn{5} Observed the following sequence of reaction :
\qfig{p2_chem_s1_q05_1}
Which of the following is/are correct?
\begin{qoptions}
\opt{A}{Product Q is \qfigi{p2_chem_s1_q05_3}}
\opt{B}{Product R is \qfigi{p2_chem_s1_q05_2}}
\opt{C}{Product P is \qfigi{p2_chem_s1_q05_4}}
\opt{D}{Reimer Tiemann reaction is involved}
\end{qoptions}

\qn{6} Compound A, C$_7$H$_{14}$ decolorizes Br$_2$ in CH$_2$Cl$_2$ and reacts with BH$_3$ in THF followed by H$_2$O$_2$ / OH$^-$ to yield compound B. When treated with KMnO$_4$, B is oxidised to a carboxylic acid C that can be resolved into enantiomers. Compound A, after ozonolysis and workup with H$_2$O$_2$, yields the same compound D as is formed by oxidation of 3-Hexanol with chromic acid. Which of the following is/are correct for compound A,B,C and D.
\begin{qoptions}
\opt{A}{Structure of compound `A' is CH$_3$--CH$_2$--CH$_2$--$\underset{\text{CH}_3}{\text{CH}}$--CH=CH$_2$}
\opt{B}{Compound `B' does not give instant turbidity with Lucas reagent.}
\opt{C}{Compound `C' is 2-Ethylpentanoic acid.}
\opt{D}{Compound `D' and 2-Hexanone are metamers.}
\end{qoptions}

\sect{SECTION-II}
\qn{1} Find the heat absorbed by an ideal gas (in kJ) when it follows the graph given below during an isothermal expansion [ln2 = 0.7] [Take : 1 litre-bar = 0.1 kJ]
\qfig{p2_chem_s2_q01_1}

\qn{2} If the mean free path is 10 cm at one bar pressure then its value in cm at 5 bar pressure, if temperature is kept constant, is

\qn{3} If bond order for CO = x

Total number of ABMO electrons in O$_2$ = y

and number of unpaired electrons in CN$^-$ = z

then calculate $\dfrac{y+z}{x}$

\qn{4} Maximum possible number of atom(s) which is/are present in the same plane of Cr(CO)$_6$ = x

Total number of unpaired electrons in [Cr(CO)$_6$] = y

Find (x + y)

\qn{5} \qfigi{p2_chem_s2_q05_1}\par
How many statements are correct?

(i) X is 1$^\circ$ alcohol

(ii) Y is 3$^\circ$ alcohol

(iii) Z is 2$^\circ$ alcohol

(iv) X, Y, Z all have phenolic OH

\qn{6} \qfigi{p2_chem_s2_q06_1}\par
x : Degree of amine in final product

y : Number of carbon atoms in Main chain of IUPAC name of final product.

Report your answer as (x $\times$ y) + 10?

\sect{SECTION-III}
\qn{1} For a certain orbital

$\Psi(r)=\dfrac{1}{16\sqrt4\,a_0^{3/2}}\,(1-\sigma)(\sigma^2-8\sigma+12)\,e^{-\sigma/2}$, $\sigma=\dfrac{2r}{na_0}$ \{$a_0$ = constant\}

The distance of nearest radial node from nucleus is $xa_0$ while farthest radial node from nucleus is $ya_0$. Find $\left(\dfrac yx\right)$.

\qn{2} The molar conductivity of a solution of a weak acid HX (0.1 M) is 100 times smaller than the molar conductivity of a solution of a weak acid HY (0.1 M). If $\lambda^0_{X^-}=\lambda^0_{Y^-}$, the difference in their pK$_a$ values, pK$_a$(HX) $-$ pK$_a$(HY), is (consider degree of ionization of both acids to be $\ll1$).

\qn{3} Find the sum of sulphide ores and chloride ores from the following given ores.

Carnallite, Argentite, Hornsilver, Rutile, Magnetite, Galena, Cinnabar, Heamatite, Zinc blende

\qn{4} Consider following reactions :\\
(1) NaH + H$_2$O $\to$\\
(2) H$_2$S + Cl$_2$ $\xrightarrow{\text{water}}$\\
(3) AgNO$_3$ + H$_2$O + H$_3$PO$_2$ $\to$\\
(4) (NH$_4$)$_2$Cr$_2$O$_{7(s)}$ $\xrightarrow{\Delta}$\\
(5) H$_2$S + SO$_{2(aq.)}$ $\to$\\
(6) XeF$_2$ + H$_2$O$_2$ $\to$

In how many of the reaction(s) at least one of the product is obtained in elemental state ?

\qfullbegin
\qn{5} Observe the following reaction sequence
\qfig{p2_chem_s3_q05_1}
\qfullsplit
Degree of unsaturation of product (Z) = X

Total number of rearrangement in reaction sequence from starting compound to Z = Y

Find X $-$ Y = ?
\qfullend

\qn{6} \qfigi{p2_chem_s3_q06_1}\par
Number of optically active isomers = x

Number of meso compounds = y

Number of geometrical isomers = z

Report your answer as x + y $-$ z = ?

\partsec{3}{MATHEMATICS}
\sect{SECTION-I}
\qn{1} Let matrix $A=[a_{ij}]$ be a square matrix of order $3\times3$, such that $1\le a_{ij}\le1000$, where $a_{ij}$ is an odd perfect square. Then which of the followings is/are always true ?
\begin{qoptions}
\opt{A}{$|C|$ is divisible by 64, where matrix $C=[c_{ij}]$ and $c_{ij}$ = cofactor of $a_{ij}$ in $|A|$.}
\opt{B}{$|A|$ is divisible by 64}
\opt{C}{Number of such matrices $A$ with all distinct elements is $^{16}C_7\cdot9!$}
\opt{D}{$|A|=0$}
\end{qoptions}

\qn{2} If $f:R-\{1\}\to R$, $f(x)=\dfrac{|x|}{(x-1)^2}$ and $[a,b]$ is an interval such that $f(a)=f(b)$, then identify the correct statement(s)
\begin{qoptions}
\opt{A}{Rolle's theorem is applicable in $[a,b]$ if $a<-1$ and $-1<b<0$.}
\opt{B}{Rolle's theorem is applicable in $[a,b]$ if $-1<a<0$ and $0<b<1$.}
\opt{C}{Rolle's theorem is applicable in $[a,b]$ if $0<a<1$ and $b>1$.}
\opt{D}{There is exactly one point on the graph of $f(x)$ where horizontal tangent can be drawn.}
\end{qoptions}

\qn{3} If $P$ and $Q$ are two square matrices of order `2' such that $P+\operatorname{adj}(Q^T)=\begin{pmatrix}2&1\\1&2\end{pmatrix}$ and $P^T-\operatorname{adj}(Q)=\begin{pmatrix}0&1\\1&0\end{pmatrix}$, then

(Note : $A^T$ denotes transpose of matrix $A$, $t_r(P)$ = sum of diagonal elements of matrix $P$)
\begin{qoptions}
\opt{A}{inverse of $Q$ exist but inverse of $P$ does not exist}
\opt{B}{$\det(Q+Q^2+Q^3)=9$}
\opt{C}{$t_r(P)=t_r(Q)$}
\opt{D}{$t_r(P)+t_r(2Q)=0$}
\end{qoptions}

\qn{4} Let $\dfrac{z_1^2}{z_2z_3}+\dfrac{z_2^2}{z_1z_3}+\dfrac{z_3^2}{z_1z_2}=3$ where $z_1$, $z_2$ and $z_3$ are there non collinear points representing the vertices of an acute angle triangle which is not equilateral. If $z_0$ is orthocentre of triangle then choose the correct option(s)-
\begin{qoptions}
\opt{A}{$|z_1-z_0|+|z_2-z_0|+|z_3-z_0|<3\left|z_1+\dfrac{z_0}{2}\right|$}
\opt{B}{$\text{amp}\left(\dfrac{z_3+\frac{z_0}{2}}{z_2+\frac{z_0}{2}}\right)=2\,\text{amp}\left(\dfrac{z_3-z_1}{z_2-z_1}\right)$, $-\pi<\text{amp}(z)\le\pi$}
\opt{C}{equation of circumcircle of $\Delta ABC$ is $\left|z+\dfrac{z_0}{2}\right|=\left|z_1+\dfrac{z_0}{2}\right|$}
\opt{D}{$\dfrac{z_1-z_0}{z_3-z_2}+\dfrac{\bar z_1-\bar z_0}{\bar z_3-\bar z_2}=0$}
\end{qoptions}

\qn{5} Let $R$ be the region bounded by $yx^4-1+y\le0$, $y\ge0$ and $0\le x\le1$. If $S_R$ be the area of region $R$, then-
\begin{qoptions}
\opt{A}{$S_R\ge\dfrac18\sqrt{\dfrac{\sqrt3}{5}}+\dfrac12$}
\opt{B}{$S_R\le\dfrac58+\dfrac38\sqrt{\dfrac{\sqrt3}{5}}$}
\opt{C}{$S_R\le\dfrac{13}{16}\sqrt{\dfrac{\sqrt3}{5}}$}
\opt{D}{$S_R\ge\dfrac34$}
\end{qoptions}

\qn{6} Let $f:R^+\to[2,\infty)$, $f(x)=2+e^{x-1}-|2x-1|+\left|e^{x-1}+|2x-1|-2\right|$ then
\begin{qoptions}
\opt{A}{$f(x)$ is bijective function}
\opt{B}{$f(x)$ is surjective function}
\opt{C}{$f(x)$ is non-differentiable at three points}
\opt{D}{Right hand derivative of $f(x)$ at $x=x_0$ is 2, where $x_0$ is point where $f(x)$ is non-differentiable $\left(x_0>\dfrac12\right)$}
\end{qoptions}

\sect{SECTION-II}
\qn{1} Let point $(1,a-1)$ lies on the diameter of circle $(x-1)^2+(y-2)^2=9$ but doesn't lie on the circle. Let $f(x)=x^2-35x-a^2$ then sum of all the integers `x' for which $f(x)$ assumes negative values, is

\qn{2} The value of $\dfrac{1}{2\pi}\displaystyle\sum_{n=1}^{20}\int_0^{\pi}\dfrac{\cos2nx-\cos6nx}{\sin^2x}\,dx$ is

\qn{3} In the given figure, each unit square is randomly coloured Red or Green. Let probability that there is no monochromatic (same colour) square of size, $2\times2$, is $p$, then value of $\dfrac{128p}{161}$ is
\qfig{p2_math_s2_q03_1}

\qn{4} Consider two straight lines $L_1:\dfrac{x-2}{-\frac12}=\dfrac{y-2}{\frac12}=\dfrac{z-4}{1}$ and $L_2:\dfrac{x-0}{\frac12}=\dfrac{y+1}{1}=\dfrac{z-3}{-1}$. Let $P$ and $Q$ be the variable points on $L_1$ and $L_2$ respectively and point $M$ divides $PQ$ in the ratio $2:3$ internally. If point $M$ moves in a plane $ax+bz=6$, then value of $a+b$ is

\qn{5} Let a parabola $y^2=\sqrt2\,x$. From a point where its directrix intersects the line $2\sqrt2\,x-\sqrt2\,y+3=0$, a normal to the parabola is drawn, meeting the parabola at $M$ (foot of normal) and $N$. $R$ is point of intersection of tangent at $M$ and chord $ON$ (produced, $O$ is origin). If angle $\angle MRN$ is $\theta$ then value of $|\tan\theta|$ is

\qn{6} Let $y=y(x)$ is the solution of the different equation $dy+e^x\tan x\,dx=y\,dx$ in $x\in\left(0,\dfrac\pi2\right)$. If $\displaystyle\lim_{x\to0^+}y(x)=0$, then area bounded by $y=y(x)$ and $f(x)=e^x(x^2-\tan x)$ is (use $e=2.71$)

\sect{SECTION-III}
\qn{1} Consider $f(x)=\dfrac{1}{1+|x-\alpha|}+\dfrac{1}{1+|x+\beta|}$. If local minimum value of $f(x)$ occurs at $x=x_0$ then $(\alpha+\beta+2)f(x_0)$ is (where $0<\alpha<\beta$)

\qn{2} Let $P(x)$ be a quadratic polynomial such that $P(1)=-1$, if $\displaystyle\int\dfrac{P(x)\,dx}{(2x-3)^2(3x-2)^2}=\dfrac15\ell n\,|f(x)|+C$ where $C$ is constant of integration and $f(x)$ is a rational function such that $\displaystyle\lim_{x\to\infty}f(x)=\dfrac43$, then value of $|f(1)|$ is

\qn{3} Let OL, OM, ON be coterminous edges of a cuboid. OL = 1 cm, OM = 2 cm, ON = 3 cm. Let $\ell$ be the shortest distance between OL and skew body diagonal, $m$ be the shortest distance between OM and skew body diagonal and $n$ be the shortest distance between ON and skew body diagonal. Then the value of $\sqrt{13\ell^2+10m^2+5n^2}$ is

\qn{4} Let $P(A\cap B)=\dfrac14$ \& $P(B)=\dfrac13$, where $P\left((A\cap B)/(A\cup\bar B)\right)=\dfrac{k}{22}$, then `k' is equal to

\qn{5} Let $m$ be the number of ways in which `$f$' number of friends visit seven different wonders such that every friend visit atleast one wonder. Let $n$ be the number of ways in which all `$f$' friends visit exactly same two wonders. If $\dfrac nm=\dfrac{21}{(127)^4}$ then value of `$f$' is.

\qn{6} Let a series of `n' ellipse $E_i:\dfrac{x^2}{a_{i+1}^2}+\dfrac{y^2}{a_i^2}=1$, where $\langle a_i\rangle$ is increasing geometric progression with common ratio `2' for $i=1,2,3,\ldots,n$. Let $C_i$ be the auxiliary circle of $E_i$. Let distance between focus and one end of minor axis of $E_n$ is $FM$ and length of that chord of $C_n$ which is tangent to $E_n$ at one end of minor axis is $L$. Then value of $\dfrac{L^2}{FM^2}$ is
\end{multicols}
\end{document}
